% ============================================================
% Apéndice A — Código y reproducibilidad
% ============================================================

\chapter{Código y reproducibilidad}
\label{ap:codigo}

Todo el código del trabajo es público en
\url{https://github.com/AgustinG-git/tfm-filamentos}. La versión que
corresponde a esta memoria está marcada con la etiqueta
\texttt{v1.0-deposito}. Este apéndice describe cómo está organizado y
cómo se reproduce cada resultado.

% ------------------------------------------------------------
\section{Organización}
% ------------------------------------------------------------

El código se divide en una biblioteca, \texttt{filseg}, con las piezas
del protocolo, y una serie de \emph{scripts} que las combinan para cada
experimento (Tablas~\ref{tab:modulos} y~\ref{tab:scripts}). Las
configuraciones de cada experimento son ficheros YAML en
\texttt{configs/}, que heredan de la baseline y cambian solo la pieza que
se estudia.

\begin{table}[ht]
\centering
\small
\begin{tabular}{@{}lp{10.5cm}@{}}
\toprule
Módulo & Contenido\\
\midrule
\texttt{contract.py} & Decisiones congeladas comunes a todas las ejecuciones (partición, resoluciones, criterio de selección, ficheros de salida). No se modifica.\\
\texttt{split.py} & Partición entrenamiento / val~$A$ / val~$B$ agrupada por fecha, generada una sola vez.\\
\texttt{gt.py} & Referencia de evaluación en el formato oficial, a partir del JSON de MAGFiLO.\\
\texttt{data.py} & Carga de imágenes a $512\times512$ y etiquetas de entrenamiento.\\
\texttt{model.py} & U-Net con codificador ResNet-34~\cite{iakubovskii2019}.\\
\texttt{losses.py} & Las pérdidas de la Tabla~\ref{tab:perdidas}: BCE, Dice, IoU, clDice y Topograph, con su calentamiento.\\
\texttt{train.py} & Una ejecución completa: entrenamiento, selección de la época en val~$A$, evaluación y registro.\\
\texttt{postproc.py} & Post-procesos: umbral y componentes conexas con filtro de área (B1, B2) filtro por persistencia (E2) y su combinación con el área (B2+P).\\
\texttt{postrun.py} & Ajuste en val~$A$ de un post-proceso sobre una ejecución ya entrenada, sin reentrenar.\\
\texttt{official\_metric.py} & Copia literal del cuaderno oficial de autoevaluación de la competición~\cite{kaggle2026}.\\
\texttt{metrics.py} & Evaluación con la métrica oficial y comprobaciones previas (formato, solapes).\\
\texttt{compare.py} & Conteos por anotación, relaciones 1:$n$ y $n$:1 y \emph{bootstrap} pareado por días.\\
\texttt{runs.py} & Identificadores, salidas de cada ejecución y tabla de registro.\\
\bottomrule
\end{tabular}
\caption{Módulos de la biblioteca \texttt{filseg}.}
\label{tab:modulos}
\end{table}

\begin{table}[ht]
\centering
\small
\begin{tabular}{@{}lp{10.5cm}@{}}
\toprule
\emph{Script} & Uso\\
\midrule
\texttt{make\_split.py} & Genera la partición (una sola vez).\\
\texttt{entrenar.py} & Lanza una ejecución a partir de una configuración.\\
\texttt{cola.sh} & Encadena varias ejecuciones en la misma GPU.\\
\texttt{postproceso.py} & Aplica B1 y B2 a las ejecuciones entrenadas (Sección~\ref{sec:res-post}).\\
\texttt{e2.py} & Aplica los post-procesos por persistencia (E2 y B2+P).\\
\texttt{ajuste\_hp.py} & Ajuste de hiperparámetros con Optuna (Sección~\ref{sec:res-hp}).\\
\texttt{resultados.py} & \emph{Bootstrap} y relaciones del Capítulo~\ref{cap:resultados}.\\
\texttt{barbas.py} & Medida de barbas y techo a $512\times512$ (Sección~\ref{sec:res-barbas}).\\
\texttt{medir\_topograph.py} & Escala del término de Topograph, para fijar $\alpha$.\\
\texttt{figura\_magfilo.py} & Figura~\ref{fig:magfilo}.\\
\bottomrule
\end{tabular}
\caption{\emph{Scripts} de los experimentos.}
\label{tab:scripts}
\end{table}

Topograph se usa a través de la implementación oficial de sus
autores~\cite{lux2025}, con su extensión compilada en C++. El
repositorio incluye además pruebas automáticas (\texttt{tests/}) que
comprueban, entre otras cosas, que la evaluación propia coincide con la
oficial y que la partición no comparte días entre subconjuntos.

% ------------------------------------------------------------
\section{Reproducir un resultado}
% ------------------------------------------------------------

Cada ejecución lee los datos de la carpeta indicada en
\texttt{FILSEG\_DATA} y escribe en \texttt{FILSEG\_OUT}. Por ejemplo, la
baseline y su post-proceso B2:

\begin{verbatim}
python scripts/entrenar.py configs/base.yaml
python scripts/postproceso.py --filtro "_base_"
\end{verbatim}

Cada ejecución crea una carpeta \texttt{runs/<run\_id>}, donde el
identificador combina la fecha, la configuración y el \emph{commit} del
código, con la configuración completa, la época elegida, las
predicciones y los \emph{logits} de validación y el envío para la
competición. Además añade una fila a \texttt{registro.csv} con sus
métricas en val~$A$ y val~$B$, y sus conteos a \texttt{conteos.csv}. Las
tablas del Capítulo~\ref{cap:resultados} salen de esos dos ficheros, que
están en el repositorio. Como el entrenamiento en GPU no es determinista,
repetir una ejecución reproduce sus cifras con una diferencia del orden
de $0{,}01$--$0{,}02$ en PQ (Sección~\ref{sec:evaluacion}).

% ------------------------------------------------------------
\section{Post-proceso por persistencia}
% ------------------------------------------------------------

Como ejemplo de cómo se traduce la teoría del
Capítulo~\ref{cap:fundamentos} a código, este es el núcleo del
post-proceso por persistencia (comentarios simplificados). Con
\texttt{area\_min} se obtiene la variante B2+P; sin él, E2. El filtro por persistencia
se apoya en los $h$-máximos: los máximos regionales de altura al menos
$\tau$ son exactamente los puntos del diagrama $H_0$ de superniveles con
persistencia $b-d\ge\tau$. Se ha comprobado, sobre imágenes sintéticas,
que la selección coincide con la que se obtiene calculando el diagrama
con \emph{union-find} y la regla del mayor
(Sección~\ref{sec:persistencia-eficiente}).

\begin{small}
\begin{verbatim}
def persistencia_512(prob, umbral, tau, delta, area_min=None):
    # X_t = {prob > umbral}, componentes 8-conexas
    xt = (prob > umbral).astype(np.uint8)
    n_t, lt = cv2.connectedComponents(xt, connectivity=8)
    # h-maximos: maximos con persistencia b - d >= tau
    hmax = h_maxima(prob, tau, footprint=np.ones((3, 3))) > 0
    keep = np.unique(lt[hmax & (xt > 0)])
    if area_min is not None:     # B2+P: area O persistencia
        areas = np.bincount(lt.ravel())
        keep = np.union1d(keep, np.flatnonzero(areas >= area_min))
    vivo = np.isin(lt, keep[keep > 0])
    # fusion por collado: unidas en X_{t - delta} -> una instancia
    _, ld = cv2.connectedComponents(
        (prob > umbral - delta).astype(np.uint8), connectivity=8)
    lab = np.where(vivo, ld, 0)
    return relabel(lab)          # etiquetas 1..n consecutivas
\end{verbatim}
\end{small}
